\section{What you can do at the end of the course}\label{appE:sec:finishing}
You do not need to remember every displayed formula in this book; what matters is what you can do. Faced with a statement, you can say exactly what it requires. Faced with a proof you do not yet follow, you can find where your understanding stops and ask for the one concept or step that would carry you past it. When an explanation arrives, from a book, a classmate or an assistant, you can test it on an example, see its role in the whole proof, and say clearly what has not yet been established.

The next unfamiliar theorem you meet will still contain mathematics you have not seen, but you now have a way in: identify the structures involved, lay out which steps depend on which, learn the missing pieces, and return to the argument. Working with an AI assistant is worthwhile when what it produces helps you understand and gives you something you can check, not when it merely fills a page with text that looks like mathematics.
